\chapter{Sums of Squares}

\begin{center}\itshape
$2026=1^2+45^2$, $2027$ is not a sum of two squares, and every integer is a sum of four.\\
Which integers are sums of two squares is decided by a single congruence on their prime factors.
\end{center}

\section{Two squares}

The \emph{Brahmagupta--Fibonacci identity} shows that sums of two squares are closed under products:
\[
(a^2+b^2)(c^2+d^2)=(ac-bd)^2+(ad+bc)^2,
\]
which is $|z|^2|w|^2=|zw|^2$ for $z=a+ib$, $w=c+id$.

\begin{lemma}[Thue]\label{lem:thue}
Let $p$ be prime and $p\nmid a$. There are integers $x,y$ with $0<|x|,|y|<\sqrt p$ and $ax\equiv y\pmod p$.
\end{lemma}
\begin{proof}
Let $m=\lfloor\sqrt p\rfloor$. The $(m+1)^2>p$ numbers $u-av$ with $0\le u,v\le m$ cannot all be distinct modulo $p$ (pigeonhole), so $u-av\equiv u'-av'$ for distinct pairs. Put $x=v-v'$ and $y=u-u'$: then $ax\equiv y$, $|x|,|y|\le m<\sqrt p$ ($p$ is not a square), and $x\ne0$ (else $y\equiv0$ forces $y=0$ and the pairs coincide), and then $y\ne0$ as $p\nmid ax$.
\end{proof}

\begin{think}
Which primes below $50$ are sums of two squares? Write them down, then reduce them modulo $4$. Why can a prime $\equiv3\pmod4$ never work? The other direction is the real theorem: what would you need to build $x$ and $y$ from $p$?
\end{think}

\begin{theorem}[Fermat, 1640; Euler, 1749]\label{thm:twosq}
An odd prime $p$ is a sum of two squares if and only if $p\equiv1\pmod4$.
\end{theorem}
\begin{proof}
Squares are $0$ or $1$ modulo $4$, so a sum of two squares is never $\equiv3\pmod4$. Conversely, if $p\equiv1\pmod4$, $-1\equiv a^2$ for some $a$ (Chapter~6). Take $x,y$ from Thue's lemma: $y^2\equiv a^2x^2\equiv-x^2$, so $p\mid x^2+y^2$, and $0<x^2+y^2<2p$. Hence $x^2+y^2=p$.
\end{proof}

\begin{theorem}\label{thm:twosqn}
An integer $n\ge1$ is a sum of two squares if and only if every prime $q\equiv3\pmod4$ appears in $n$ with an even exponent.
\end{theorem}
\begin{proof}
If the condition holds, $n=m^2\cdot2^a\cdot\prod p_i$ with $p_i\equiv1\pmod4$; each factor is a sum of two squares ($m^2=m^2+0^2$, $2=1^2+1^2$), and so is the product. Conversely, let $n=x^2+y^2$ and $q\equiv3\pmod4$ divide $n$. If $q\nmid x$, then $\left(yx^{-1}\right)^2\equiv-1\pmod q$, impossible. So $q\mid x$, $q\mid y$, $q^2\mid n$, and $\frac n{q^2}=\left(\frac xq\right)^2+\left(\frac yq\right)^2$. Induction on $v_q(n)$ shows $v_q(n)$ is even.
\end{proof}

Counting representations is subtler. Let $r_2(n)$ be the number of $(x,y)\in\Z^2$ with $x^2+y^2=n$ (signs and order count). Then
\[
r_2(n)=4\sum_{d\mid n}\chi_4(d),\qquad\chi_4(d)=\begin{cases}0&d\text{ even},\\1&d\equiv1\pmod4,\\-1&d\equiv3\pmod4,\end{cases}
\]
a formula of Jacobi. \emph{We do not use it here; it is proved in Chapter~10 with Gaussian integers.} For instance $r_2(25)=4(1+1+1)=12$: $(\pm5,0),(0,\pm5),(\pm3,\pm4),(\pm4,\pm3)$.

\section{Four squares}

\begin{think}
Find the integers $n\le30$ that are not sums of three squares. What do they have in common modulo $8$? Do any of them need more than four squares?
\end{think}

\begin{theorem}[Lagrange, 1770]\label{thm:foursq}
Every positive integer is a sum of four squares.
\end{theorem}
\begin{proof}
\emph{Euler's identity.} One checks by expansion (or with quaternion norms, Tome~2) that
\[
\left(\sum_{i=1}^4a_i^2\right)\left(\sum_{i=1}^4b_i^2\right)=z_1^2+z_2^2+z_3^2+z_4^2,
\]
with $z_1=a_1b_1+a_2b_2+a_3b_3+a_4b_4$, $z_2=a_1b_2-a_2b_1+a_3b_4-a_4b_3$, $z_3=a_1b_3-a_3b_1+a_4b_2-a_2b_4$, $z_4=a_1b_4-a_4b_1+a_2b_3-a_3b_2$. So it suffices to treat primes, and $2=1+1+0+0$.

\emph{A first multiple.} Let $p$ be an odd prime. The $\frac{p+1}2$ numbers $x^2$, $0\le x\le\frac{p-1}2$, are distinct modulo $p$, and so are the $\frac{p+1}2$ numbers $-1-y^2$. Together they are $p+1$ numbers, so some $x^2\equiv-1-y^2$: $x^2+y^2+1+0^2=mp$ with $0<mp\le2\left(\frac{p-1}2\right)^2+1<p^2$, so $1\le m<p$.

\emph{Descent.} Let $m$ be the least positive integer with $mp=x_1^2+x_2^2+x_3^2+x_4^2$; we know $m<p$. If $m$ is even, an even number of the $x_i$ are odd; pair them so that $x_1\equiv x_2$, $x_3\equiv x_4\pmod2$, and $\frac m2p=\left(\frac{x_1+x_2}2\right)^2+\left(\frac{x_1-x_2}2\right)^2+\left(\frac{x_3+x_4}2\right)^2+\left(\frac{x_3-x_4}2\right)^2$ contradicts minimality. So $m$ is odd. Suppose $m>1$. Choose $y_i\equiv x_i\pmod m$ with $|y_i|\le\frac{m-1}2$. Then $\sum y_i^2\equiv\sum x_i^2\equiv0\pmod m$, say $\sum y_i^2=mr$, with $r\le\frac{4}{m}\left(\frac{m-1}2\right)^2<m$. And $r>0$: otherwise all $x_i\equiv0\pmod m$, so $m^2\mid mp$, $m\mid p$, impossible for $1<m<p$.
Now apply Euler's identity to $(mp)(mr)$ with $a=x$, $b=y$. Since $y_i\equiv x_i\pmod m$: $z_1\equiv\sum x_i^2\equiv0$, and each of $z_2,z_3,z_4$ is $\equiv0$ because its terms cancel in pairs modulo $m$ (e.g.\ $x_1y_2-x_2y_1\equiv x_1x_2-x_2x_1$). So $rp=\sum\left(\frac{z_i}m\right)^2$ with $0<r<m$, contradicting minimality. Hence $m=1$.
\end{proof}

\begin{remark}
Three squares are harder: \emph{Legendre and Gauss} proved that $n$ is a sum of three squares if and only if $n$ is not of the form $4^a(8b+7)$. The ``only if'' direction is Exercise~\ref{ex:three}; the converse needs either Gauss' theory of ternary quadratic forms or Dirichlet's theorem on primes in progressions, and is proved in Tome~12 (with the Hasse--Minkowski theorem). We only use it in Exercise~\ref{ex:eureka}, where it is explicitly admitted.
\end{remark}

\section{Lattice points in a disc}

\begin{think}
How many points of $\Z^2$ lie in the disc of radius $10$ centred at $0$? Guess first, using only the area of the disc. (The answer is $317$, against $100\pi\approx314.2$.) Why should the area be a good guess, and how large could the error be?
\end{think}

\begin{theorem}[Gauss]\label{thm:circle}
Let $N(R)=\#\{(x,y)\in\Z^2:x^2+y^2\le R^2\}$. Then $\left|N(R)-\pi R^2\right|\le2\sqrt2\,\pi R+2\pi$ for $R\ge1$. In particular $N(R)\sim\pi R^2$.
\end{theorem}
\begin{proof}
Attach to each counted point $(x,y)$ the unit square $[x,x+1[\times[y,y+1[$. These squares are disjoint, have total area $N(R)$, and lie in the disc of radius $R+\sqrt2$ (a point of a square is within $\sqrt2$ of its corner). They cover the disc of radius $R-\sqrt2$: a point $P$ of that disc lies in the square of the lattice point $(\lfloor P_x\rfloor,\lfloor P_y\rfloor)$, which is within $\sqrt2$ of $P$, hence in the disc of radius $R$. So $\pi(R-\sqrt2)^2\le N(R)\le\pi(R+\sqrt2)^2$, and expand.
\end{proof}

\begin{figure}[ht]
\centering
\fig{ch07-disc}
\caption{The $N(6)=113$ lattice points in the disc of radius $6$, against $36\pi\approx113.1$. Each point $(x,y)$ owns the unit square $[x,x+1[\times[y,y+1[$. The squares cover the disc of radius $R-\sqrt2$ and stay inside the disc of radius $R+\sqrt2$ (dashed), which squeezes $N(R)$ between the two areas.}
\label{fig:disc}
\end{figure}

Since $N(R)=\sum_{n\le R^2}r_2(n)$, the \emph{average} of $r_2(n)$ is $\pi$. The true size of the error is the \emph{Gauss circle problem}: it is believed to be $O(R^{1/2+\varepsilon})$, and the best bound known today is about $O(R^{0.63})$.

\section{Exercises}

\begin{exercise}[$\star$]\label{ex:three}
Show that an integer of the form $4^a(8b+7)$ is not a sum of three squares. \emph{Look at squares modulo $8$.}
\end{exercise}

\begin{exercise}[$\star$]\label{ex:2026}
Write $1013$ (a prime) and then $2026$ as sums of two squares. Explain why $2027$ is not a sum of two squares, and find all $n\in\{2021,\dots,2030\}$ that are.
\end{exercise}

\begin{exercise}[$\star\star$]\label{ex:primitive2}
Show that if $n=x^2+y^2$ with $\gcd(x,y)=1$, then $4\nmid n$ and $n$ has no prime factor $\equiv3\pmod4$.
\end{exercise}

\begin{exercise}[$\star\star$]\label{ex:x2y2}
Show that a prime $p\equiv1$ or $3\pmod8$ is of the form $x^2+2y^2$. \emph{Use Thue's lemma with $a^2\equiv-2$ (Chapter~6), then deal with the case $x^2+2y^2=2p$.}
\end{exercise}

\begin{exercise}[$\star\star$]\label{ex:x3y2}
Show that a prime $p\equiv1\pmod3$ is of the form $x^2+3y^2$.
\end{exercise}

\begin{exercise}[$\star\star$]\label{ex:eureka}
\emph{Gauss' ``E\upshape{}YPHKA'' (1796).} Admitting the three-square theorem, show that every $8n+3$ is a sum of three odd squares, and deduce that every $n\in\N$ is a sum of three triangular numbers $\frac{k(k+1)}2$.
\end{exercise}

\begin{exercise}[$\star\star\star$]\label{ex:zagier}
\emph{Zagier's one-sentence proof (1990).} Let $p\equiv1\pmod4$ be prime and $S=\{(x,y,z)\in(\N^*)^3:x^2+4yz=p\}$.
\begin{enumerate}[label=(\alph*)]
\item Show that $S$ is finite and that $x\ne y-z$ and $x\ne2y$ for $(x,y,z)\in S$.
\item Show that the following map is an involution of $S$ (it sends $S$ to $S$ and $f\circ f=\operatorname{id}$):
\[
f(x,y,z)=\begin{cases}(x+2z,\ z,\ y-x-z)&\text{if }x<y-z,\\(2y-x,\ y,\ x-y+z)&\text{if }y-z<x<2y,\\(x-2y,\ x-y+z,\ y)&\text{if }x>2y.\end{cases}
\]
\item Show that $f$ has exactly one fixed point. Deduce that $|S|$ is odd.
\item Deduce that the involution $(x,y,z)\mapsto(x,z,y)$ has a fixed point, and that $p$ is a sum of two squares.
\end{enumerate}
\end{exercise}

\begin{beauty}
Two squares, three squares, four squares: the answers are, respectively, a condition on primes $\equiv3\pmod4$, a condition modulo $8$, and ``always''. The two-square theorem becomes transparent in the Gaussian integers (Chapter~10), where $p=x^2+y^2$ means $p=(x+iy)(x-iy)$; the four-square theorem in the quaternions (Tome~2). The study of which primes are $x^2+ny^2$ grows into class field theory, and Zagier's proof is a reminder that sometimes a clever involution is all one needs.
\end{beauty}

\section{Hints}

\begin{hint}{ex:three}
Squares are $0,1,4\pmod8$. Which residues modulo $8$ can a sum of three squares take? For the factor $4^a$, show that if $4\mid x^2+y^2+z^2$ then $x,y,z$ are all even.
\end{hint}
\begin{hint}{ex:2026}
Factor each number and apply Theorem~\ref{thm:twosqn}. To find the squares, write each prime $\equiv1\pmod4$ as a sum of two squares and multiply with the Brahmagupta--Fibonacci identity.
\end{hint}
\begin{hint}{ex:primitive2}
Reread the proof of Theorem~\ref{thm:twosqn}: a prime $q\equiv3\pmod4$ dividing $x^2+y^2$ divides both $x$ and $y$. For $4$, look at $x^2+y^2$ modulo $4$ when $x$ and $y$ are not both even.
\end{hint}
\begin{hint}{ex:x2y2}
With $a^2\equiv-2$ and Thue's $x,y$, you get $p\mid y^2+2x^2<3p$. So $y^2+2x^2$ is $p$ or $2p$. In the second case $y$ is even: write $y=2y'$.
\end{hint}
\begin{hint}{ex:x3y2}
Now $y^2+3x^2<4p$, so it equals $p$, $2p$ or $3p$. Show $2p$ is impossible modulo $4$ or $8$, and in the case $3p$ note that $3\mid y$.
\end{hint}
\begin{hint}{ex:eureka}
$8T_k+1=(2k+1)^2$. If $8n+3=a^2+b^2+c^2$, look modulo $8$ to see that $a,b,c$ must all be odd.
\end{hint}
\begin{hint}{ex:zagier}
(a) $x\ne y-z$ and $x\ne2y$ because $p$ is prime and not a square. (c) A fixed point must lie in the middle case with $x=y$; then $x(x+4z)=p$. (d) An involution of a finite set has as many fixed points as the size of the set, modulo $2$.
\end{hint}

\section{Solutions}

\begin{solution}{ex:three}
\idea A statement of the form ``never a sum of three squares'' is usually proved by a congruence: find a modulus where squares take few values. Modulo $4$, squares are $0,1$; modulo $8$, odd squares are all $1$ (since $(2k+1)^2=4k(k+1)+1$ and $k(k+1)$ is even). So try modulo $8$, and handle the factor $4^a$ by descent.
\writeup Squares are $0,1,4\pmod8$, and a sum of three of them takes the values $0,1,2,3,4,5,6$ but never $7$. If $4\mid x^2+y^2+z^2$, then modulo $4$ (squares $0,1$) the only way three squares sum to $0$ is $0+0+0$: all three are even, and $\frac n4=\left(\frac x2\right)^2+\left(\frac y2\right)^2+\left(\frac z2\right)^2$. So if $4^a(8b+7)$ were a sum of three squares, dividing by $4$ $a$ times would make $8b+7$ one: impossible.
\end{solution}

\begin{solution}{ex:2026}
\idea Theorem~\ref{thm:twosqn} turns ``is $n$ a sum of two squares?'' into ``factor $n$ and look at primes $\equiv3\pmod4$''. To actually find the squares, decompose the prime factors $\equiv1\pmod4$ and multiply with the Brahmagupta--Fibonacci identity. For $1013$, search $x^2=1013-y^2$ near $\sqrt{1013/2}\approx22.5$.
\writeup $1013-23^2=484=22^2$, so $1013=22^2+23^2$. With $2=1^2+1^2$ and the identity ($a=b=1$, $c=22$, $d=23$): $2026=(22-23)^2+(22+23)^2=1^2+45^2$. $2027$ is prime and $\equiv3\pmod4$, so it is not a sum of two squares. In $\{2021,\dots,2030\}$: $2021=43\cdot47$, $2022=2\cdot3\cdot337$, $2023=7\cdot17^2$, $2024=2^3\cdot11\cdot23$, $2028=2^2\cdot3\cdot13^2$ and $2030=2\cdot5\cdot7\cdot29$ each contain a prime $\equiv3\pmod4$ to an odd power: no. $2025=45^2+0^2$, $2026$, and $2029$ (a prime $\equiv1\pmod4$, equal to $2^2+45^2$) are sums of two squares.
\end{solution}

\begin{solution}{ex:primitive2}
\idea Both claims come from the proof of Theorem~\ref{thm:twosqn}, which shows that a prime $q\equiv3\pmod4$ dividing $x^2+y^2$ divides both $x$ and $y$; coprimality forbids that. For $4$, look modulo $4$.
\writeup If $4\mid x^2+y^2$, then since squares are $0,1\pmod4$, both squares are $0\pmod4$, so $x,y$ are even, contradicting $\gcd(x,y)=1$. If a prime $q\equiv3\pmod4$ divided $n$ but not $x$, then $\left(yx^{-1}\right)^2\equiv-1\pmod q$, impossible since $-1$ is not a square modulo $q$ (Chapter~6). So $q\mid x$, then $q\mid y^2$, so $q\mid y$: contradiction.
\end{solution}

\begin{solution}{ex:x2y2}
\idea Copy the proof of Theorem~\ref{thm:twosq}: that proof needed (1) a square root $a$ of $-1$ modulo $p$, and (2) Thue's lemma to make $x,y$ small. Here the form is $x^2+2y^2$, so we need $a^2\equiv-2$, which Chapter~6 provides exactly for $p\equiv1,3\pmod8$. Thue then gives a multiple of $p$ that is small, and we finish by checking the few possible multiples.
\writeup Let $a^2\equiv-2\pmod p$ and take $x,y$ from Thue's lemma: $y\equiv ax$, so $y^2+2x^2\equiv(a^2+2)x^2\equiv0\pmod p$, with $0<y^2+2x^2<p+2p=3p$. So $y^2+2x^2\in\{p,2p\}$. If it is $p$, done. If it is $2p$, then $y^2$ is even, so $y=2y'$, and $4y'^2+2x^2=2p$ gives $x^2+2y'^2=p$.
\end{solution}

\begin{solution}{ex:x3y2}
\idea Same strategy with $-3$: by Chapter~6 (Exercise~6.1), $-3$ is a square modulo $p$ exactly when $p\equiv1\pmod3$. Thue's lemma now bounds the multiple by $4p$, so there are three cases to examine, and the cases that do not directly give $p$ must be either converted ($3p$) or ruled out by a parity argument ($2p$).
\writeup Let $a^2\equiv-3$ and take $x,y$ from Thue's lemma: $y^2+3x^2\equiv0\pmod p$ with $0<y^2+3x^2<4p$, so it equals $p$, $2p$ or $3p$.
If it is $3p$: then $3\mid y^2$, so $y=3y'$, and $9y'^2+3x^2=3p$ gives $x^2+3y'^2=p$.
If it is $2p$: then $y^2+3x^2$ is even, so $x\equiv y\pmod2$. If both are even, $4\mid2p$; if both are odd, $y^2+3x^2\equiv1+3\equiv0\pmod4$, so again $4\mid2p$. Both contradict $p$ odd. So the value is $p$ or $3p$, and in both cases $p=X^2+3Y^2$.
\end{solution}

\begin{solution}{ex:eureka}
\idea Triangular numbers are linked to odd squares by $8T_k+1=(2k+1)^2$. So ``$n$ is a sum of three triangular numbers'' is the same as ``$8n+3$ is a sum of three \emph{odd} squares''. The three-square theorem gives three squares; a look modulo $8$ forces them to be odd.
\writeup $8n+3$ is not of the form $4^a(8b+7)$ (it is odd and $\equiv3\pmod8$), so $8n+3=a^2+b^2+c^2$. Squares are $0,1,4\pmod8$, and the only way to reach $3$ with three of them is $1+1+1$, so $a,b,c$ are odd: $a=2k+1$, $b=2l+1$, $c=2m+1$. Since $(2k+1)^2=8T_k+1$ with $T_k=\frac{k(k+1)}2$, we get $8n+3=8(T_k+T_l+T_m)+3$, so $n=T_k+T_l+T_m$.
\end{solution}

\begin{solution}{ex:zagier}
\idea The trick is a parity argument: an involution on a finite set has a number of fixed points with the same parity as the set. The set $S$ is designed so that one involution, $(x,y,z)\mapsto(x,z,y)$, has as fixed points exactly the decompositions $p=x^2+(2y)^2$; Zagier's second involution $f$ shows $|S|$ is odd. So the work is checking $f$ (routine algebra) and finding its fixed points.
\writeup (a) $x^2<p$ and $yz<\frac p4$ leave finitely many triples. If $x=y-z$ then $p=(y-z)^2+4yz=(y+z)^2$, and if $x=2y$ then $p=4y(y+z)$; neither is prime. So the three cases of $f$ cover $S$.

(b) In each case the image has positive coordinates (this is what the case conditions say: $y-x-z>0$ in the first, $2y-x>0$ in the second, $x-2y>0$ in the third) and $x'^2+4y'z'=x^2+4yz$; for instance in the first case $(x+2z)^2+4z(y-x-z)=x^2+4xz+4z^2+4yz-4xz-4z^2=x^2+4yz$, and the other two are similar. The first case lands in the third ($x'=x+2z>2z=2y'$), the third in the first ($x'=x-2y<x-y+z-y=y'-z'$), and the second in the second; substituting shows $f(f(x,y,z))=(x,y,z)$ in each case.

(c) In the first and third cases $f$ changes the case, so there is no fixed point. In the second, $f(x,y,z)=(x,y,z)$ means $2y-x=x$, i.e.\ $x=y$. Then $p=x^2+4xz=x(x+4z)$ with $p$ prime forces $x=1$, so $y=1$ and $z=\frac{p-1}4$. So $f$ has exactly one fixed point. The other elements of $S$ come in pairs $\{s,f(s)\}$, hence $|S|$ is odd.

(d) $g(x,y,z)=(x,z,y)$ is also an involution of $S$. Its non-fixed points come in pairs, and $|S|$ is odd, so it has a fixed point: $y=z$, and $p=x^2+4y^2=x^2+(2y)^2$.
\end{solution}

\section*{Sources and further reading}

\begin{itemize}
\item E.~Grosswald, \emph{Representations of Integers as Sums of Squares}, Springer, 1985. Everything about two, three, four and more squares, with $r_k(n)$ formulas.
\item D.~Zagier, A one-sentence proof that every prime $p\equiv1\pmod4$ is a sum of two squares, \emph{American Mathematical Monthly} 97 (1990), 144. The proof of Exercise~\ref{ex:zagier}.
\item D.~A.~Cox, \emph{Primes of the Form $x^2+ny^2$}, 3rd ed., AMS Chelsea, 2022, \S\S1--3. Fermat's questions about $x^2+2y^2$ and $x^2+3y^2$ (Exercises~\ref{ex:x2y2} and~\ref{ex:x3y2}) and how they lead to class field theory.
\item M.~N.~Huxley, Exponential sums and lattice points~III, \emph{Proceedings of the London Mathematical Society} 87 (2003), 591--609. The bound $O(R^{131/208+\varepsilon})$, about $R^{0.63}$, for the circle problem.
\item J.~Bourgain and N.~Watt, Mean square of zeta function, circle problem and divisor problem revisited, preprint, arXiv:1709.04340, 2017. A more recent improvement, to $R^{517/824+\varepsilon}\approx R^{0.6275}$, and a glimpse of current research.
\end{itemize}
